% Part: sets-functions-relations
% Chapter: size-of-sets-complete

\documentclass[../../../include/open-logic-chapter]{subfiles}

\begin{document}

\olchapter{sfr}{siz}{Kümelerin Büyüklüğü}

\begin{editorial}
Bu bölüm numaralandırmaları, sayılabilirliği ve sayılamazlığı ele alır.
Bazı kesimlerin iki sürümü vardır: daha temel sürüm, numaralandırmaları
listeler veya $\PosInt$'tan örten fonksiyonlar olarak ele alır; daha soyut
sürüm ise numaralandırmaları $\Nat$ ile bire bir örten eşlemeler olarak
tanımlar.
\end{editorial}

\olimport{introduction}

\olimport{enumerability}

\olimport{zig-zag}

\olimport{pairing}

\olimport{pairing-alt}

\olimport{non-enumerability}

\olimport{reduction}

\olimport{equinumerous-sets}

\olimport{comparing-size}

\olimport{schroder-bernstein}

\begin{editorial}
Aşağıdaki \olref[sfr][siz][enm-alt]{sec},
\olref[sfr][siz][nen-alt]{sec} ve \olref[sfr][siz][red-alt]{sec}, Tim
Button'ın \emph{Open Set Theory} metninde kullanılmak üzere hazırladığı,
sırasıyla \olref[sfr][siz][enm]{sec}, \olref[sfr][siz][nen]{sec} ve
\olref[sfr][siz][red]{sec} kesimlerinin alternatif sürümleridir. Bunlar biraz
daha ileri düzeydedir ve küme teorisi bağlamına daha uygun farklı bir
sayılabilirlik tanımı kullanır (yani listelenebilirlik veya $\PosInt$'tan örten
bir fonksiyonun görüntü kümesi olmak yerine $\Nat$ ya da onun bir başlangıç
parçasıyla bire bir örten eşlenebilirlik).
\end{editorial}

\olimport{enumerability-alt}
\olimport{non-enumerability-alt}
\olimport{reduction-alt}

\OLEndChapterHook

\end{document}
